\section{The Schema, as a String}
\label{sec:ch06-the-schema-as-a-string}

\index{service@\code{\_service}!what it answers|(}%
The first of the two new fields returns your schema. Ask for it the way you
would ask for anything:

\begin{minted}{text}
POST /graphql HTTP/1.1
Host: localhost:5001
Content-Type: application/json

{"query":"{ _service { sdl } }"}
\end{minted}

\index{service@\code{\_service}!returns the exported document}%
What comes back is the whole schema inside one JSON string. Undo the escaping,
normalize the line endings, and you have the file \code{schema export} writes,
byte for byte. It is the same document rather than a federation-only view of
one, descriptions and \code{@cost} directives and all.

\index{service@\code{\_service}!costs no statement}%
It also costs nothing to serve. Zero statements: the schema is in memory and no
table is touched. That matters more than it sounds, because of who asks and how
often.

\subsection{Who Reads It}

\index{service@\code{\_service}!is for a router at startup}%
Apollo's specification puts \code{\_service} under a heading that answers the
question: \enquote{Enhanced introspection with \code{Query.\_service}}. The
reader it has in mind is a router working out a supergraph for
itself~\autocite{apollosubgraphspec}:

\begin{quote}
Some federated graph routers can compose their supergraph schema dynamically at
runtime. To do so, a graph router first executes the following enhanced
introspection query on each of its subgraphs to obtain all subgraph schemas:
\end{quote}

\index{composition!dynamic composition is discouraged}%
And then, unusually, argues against the use case it has just described. Dynamic
composition \enquote{can cause unexpected downtime if composition fails on
router startup}, so Apollo \enquote{strongly recommends against} it, while
holding that \enquote{supporting this use case is still a requirement for
subgraph libraries}~\autocite{apollosubgraphspec}. Chapter~\ref{ch:composition}
composes ahead of time with a command-line tool for exactly that reason, and
this field is still required to exist.

\index{sdl@\code{sdl}!what the string must contain}%
What the string must contain is one clause with a condition on it. It
\enquote{must include all uses of all federation-specific directives, such as
\code{@key}}, and whether the library's own additions may appear depends on
whether that library still supports Federation~1: one that does must leave them
out, and one that supports only Federation~2 \enquote{can include these
definitions}~\autocite{apollosubgraphspec}. So how much of the document you get
back is a property of the library rather than of the specification, and reading
it off a running service is the only way to know. This one includes everything.

\index{sdl@\code{sdl}!the description disagrees with the field}%
Which produces a contradiction a reader who introspects this type will meet.
The generated description on \code{\_Service} says the returned document
\enquote{does not include the additions of the federation spec}. It does.
\code{\_service}, \code{\_entities}, \code{\_Any}, \code{\_Entity} and the
directive definitions are all in the string this service hands back. I take the
string over the sentence describing it, here and generally: the description was
written once and the string is produced on every request.

\subsection{The Field Nobody Meant You to Have}

\index{service@\code{\_service}!is not hidden}%
Beside the sentences about routers, the specification adds a note. The field
\enquote{is not included in the composed supergraph schema}, and \enquote{for
security reasons, it's intended solely for use by the graph
router}~\autocite{apollosubgraphspec}.

\index{service@\code{\_service}!is visible to introspection}%
Read that against what you can see from outside. Ask this service for the fields
on \code{Query} and it lists \code{node}, \code{nodes}, \code{sessions},
\code{\_service} and \code{\_entities}. No authorization, no separate port, no
header. \code{sessionById} is missing from that list and the two underscore
fields are not. Read that pairing again: the field this book deprecated in
chapter~\ref{ch:schema-design} is hidden from a default introspection query, and
the two fields the specification calls internal are not.

\index{service@\code{\_service}!intent is not a control}%
\enquote{Intended solely for} is a statement about intent, and a subgraph has no
mechanism behind it. What the sentence accurately describes is the supergraph:
the composed document a client reads does not carry the field, so a client
talking to the router cannot select it. A client talking to your service can, and
in chapter~\ref{ch:enter-the-router} your service is still on a port. Whether
anything else can reach that port is your deployment's answer and not the
specification's, and chapter~\ref{ch:entities-authorization} is the chapter about
assuming otherwise.
\index{service@\code{\_service}!what it answers|)}%
